\documentclass[11pt]{article}
\usepackage{amsmath,amssymb,calc,color}
\usepackage{graphicx}
\usepackage{amsthm}%
\usepackage{natbib}
\usepackage{bm}
%\usepackage{macros}
\usepackage[latin1]{inputenc}
\makeatletter

\newdimen\breite
\setlength{\breite}{29pt}
\newdimen\hohe
\setlength{\hohe}{26pt}

\addtolength{\textwidth}{4\breite}
\addtolength{\oddsidemargin}{-2\breite}
\addtolength{\evensidemargin}{-2\breite}
\addtolength{\textheight}{4\hohe}
\addtolength{\topmargin}{-1.95\hohe}


\renewcommand\baselinestretch{1.2}%{1.82}

\def\ul{\bgroup\smallskip
\parindent-15pt
\leftskip=15pt}
\def\endul{\vskip1sp\smallskip\egroup}

\definecolor{DarkBlue}{rgb}{0,0,0.5451}

\newcommand{\answer}[1]{\par\smallskip\noindent{\color{DarkBlue}\it#1}\par\medskip}
\newcommand{\question}[1]{\par\smallskip\noindent$\leadsto$ {\color{red}\sc#1}\par\medskip}
\newcommand{\bs}{\bigskip}
\setcounter{section}1
\renewcommand\thesubsection   {{\bf\@arabic\c@subsection}}




\begin{document}
\parindent=0pt


\newpage
\begin{enumerate}
\item Page 7. Can we make a statement like: If $w \sim PG(b,c)$ then $w \sim PIG(0,\sqrt{\frac{2}{2\pi^2(k-1/2)^2 + c^2/2}})$?
\answer{PG is PIG only if $b = -\frac{3}{2}$ and $d_k = \infty$ for all $k$.}

\item Page 7. Are these just two examples of ratios of gamma functions? I don't immediately see the connection between equations 9 and 10. Am I missing something? I do think it's a good idea to connect the ratio of gamma functions to the scale mixture representation early on.
\answer{Maybe we should delete this remark.}

\item Page 8. Is this only true for $a = 1$ or does it generalize to other values?
\answer{The result $$
\frac{\Gamma(a)}{\Gamma(a+t)} e^{\psi(a)t} = E_{w|a}\left[ e^{-\frac{1}{2}t^2 w}\right]$$ holds for all $a>0$ and $t\geq 0$. $\psi(\cdot)$ is digamma function.}

\item Page 11. I think we still need a proof of the full conditional for $w_{mk}$ here.
\answer{
By Equation 19, when condition on all other variables, the density of $w_{mk}$ is proportional to 
$$
e^{-(n_{mk}+\alpha_k-1)^2w_{mk}}p(w_{mk})
$$
where $p(\cdot)$ is the density of $PIG(\bm{d},0)$. Then compare the above expression with Equation 6, we see that the conditional distribution of $w_{mk}$ is $PIG(\bm{d}, \sqrt{2(n_{mk}+\alpha_k-1)^2})$.}

\end{enumerate}


\end{document}